\section{From source contours to covering words}
\label{sec-assembly}

The numerical recurrence used later concerns finite lists of rows and
circles. We describe those lists and prove that an extension may use any
positive number of new letters. The one-letter transport lemmas in the
supplied assembly library apply repeatedly to give this construction.

\subsection{Rows, closed contours, and insertion}

Write $\operatorname{rot}(b,j)$ for the word obtained by rotating a nonempty
word $b$ to the left through $j$ positions. A \emph{cyclic class} consists
of all rotations of a word with distinct letters. If $x$ is absent from
$b$, write $\operatorname{ins}(b,x,j)$ for insertion of $x$ immediately
before position $j$, where positions start at zero and position $|b|$
means insertion at the end. Thus
$\operatorname{ins}(abc,x,1)=axbc$ and
$\operatorname{ins}(abc,x,3)=abcx$.

A \emph{row} is a triple $(b,s,v)$, where $b$ has distinct letters, the
letter $s$ does not occur in $b$, and $1\le v\le |b|$. The word $b$ is its
\emph{base}, the letter $s$ is its \emph{satellite}, and its \emph{charge}
is $|b|-v$. Its head and tail are the following literal words:
\[
 \operatorname{head}(b,s,v)=b[0:|b|-2],\qquad
 \operatorname{tail}(b,s,v)
   =\operatorname{rot}(b,v+1)[0:|b|-2].
\]
Here $b[i:j]$ means the consecutive letters in positions $i$ through
$j-1$. A row is \emph{full} if its charge is zero, and \emph{short} if
its charge is two. Except during the final completion, all rows below are
full or short. All rows in a given contour have the same satellite and the
same base alphabet.

Two rows are \emph{compatible} if the first tail equals the second head
as a word. A \emph{closed contour} is a nonempty finite list of rows in
which each row is compatible with the next, including the last and first
rows. Repeated occurrences remain separate entries of the list. For
example, the full rows with bases $abcd$, $bcad$, and $cabd$, all with one
fixed satellite outside $\{a,b,c,d\}$, form a closed contour. Their heads
are $ab,bc,ca$, respectively, and their tails are $bc,ca,ab$.

For a base of length $h+1$, an \emph{insertion port} is an integer in
$\{0,\ldots,h-1\}$. These are the $h$ insertion positions in a head of
length $h-1$. Across a full row the port changes from $p$ to $p-1$ modulo
$h$, and across a short row it changes from $p$ to $p+1$ modulo $h$.
Consequently the port displacement of a row is its charge minus one.
If a contour has $N$ rows and total charge $q$, its displacement is
$\Delta=q-N$.

We recall the literal transport operation from the supplied assembly
library~\cite{pantone-core}. Let $z$ be a fresh ordinary letter. A full
row $(b,s,h+1)$ at port zero is replaced by the two full rows with bases
$\operatorname{ins}(b,z,0)$ and $\operatorname{ins}(b,z,h)$.
At a port $1\le p<h$ it is replaced by the single full row with base
$\operatorname{ins}(b,z,p)$. A short row at a port $p<h-1$ is replaced
by the single short row with base $\operatorname{ins}(b,z,p)$.
At port $h-1$ it is replaced by two short rows with bases
$\operatorname{ins}(b,z,h-1)$ and
$\operatorname{rot}(b,h)z$. The satellite is unchanged in all four cases.
A transported contour is formed by concatenating these local paths and
propagating the port between successive source rows.

Each local path starts at the old head with $z$ inserted at the input
port and ends at the old tail with $z$ inserted at the output port.
This follows directly from the displayed definitions of head and tail.
For example, when $b=aUb$ is a full base and the input port is zero,
the two new bases are $zaUb$ and $aUzb$. Their common interface is the
literal word $aU$. These endpoint equations explain why compatible
source rows give compatible transported paths.

Across all $h$ input ports, each source row gives exactly $h+1$ new row
occurrences. All those occurrences have its charge. Their base cyclic
classes are precisely the classes obtained by inserting $z$ in all cyclic
gaps of the source base. The endpoint, inventory, and coverage statements
are the local transport propositions imported from the supplied
library~\cite{pantone-core}. The formulas above specify the operation to
which those propositions apply.

\begin{lemma}\label{lem-assembly-orbits}
Let a closed contour have bases of length $h+1$, where $h\ge2$, and put
$g=\gcd(h,|\Delta|)$ and $\ell=h/g$. For each $0\le e<g$, transport
$\ell$ consecutive copies of the contour from initial port $e$.
The resulting $g$ contours are closed. Their displacements are divisible
by $h+1$, and their combined row and charge counts are $(h+1)N$ and
$(h+1)q$, respectively. Every source-row and source-port pair occurs once
in their transport.
\end{lemma}

\begin{proof}
The ports at the beginnings of the successive source laps in component
$e$ are $e+t\Delta$ modulo $h$, for $0\le t<\ell$. Consider the map
\[
 (e,t)\longmapsto e+t\Delta\pmod h,
 \qquad 0\le e<g,\quad 0\le t<h/g.
\]
Equality of two outputs first gives equality of $e$ modulo $g$, and thus
equality of $e$ in its specified range. Division of the remaining
congruence by $g$ then gives
$(t-t')\Delta/g\equiv0\pmod{h/g}$. The integers $\Delta/g$ and $h/g$
are coprime, so $t=t'$ in the specified range. The domain and codomain
both have $h$ elements, so the map is a bijection. If $\Delta=0$, this
argument reduces to $g=h$, $\ell=1$, and the identity enumeration of
the ports.

After $\ell$ laps the source port returns to $e$, since $h$ divides
$\ell\Delta$. The literal endpoint equations therefore close the
transported path. The winding proposition for transport gives displacement
$(h+1)\Delta/g$ for that path, which is divisible by $h+1$. This is the
repeated-contour case of the imported winding theorem; its divisibility
hypothesis is exactly $h\mid\ell\Delta$.

At any fixed source-row occurrence, the prefix port change is a translation
of the set of lap-start ports. The bijection above therefore shows that
all its input ports occur once across the $g$ components. The local
inventory now gives $h+1$ output rows, each with the source charge, for
every source-row occurrence. Summing proves both count formulas.
\end{proof}

\subsection{Explicit circles meeting every completion port}

We reserve a second fresh letter $f$ for final completion. For precision,
we describe the completion operation used in the imported joining theorem.
For a row $(b,s,v)$, its lift at position $j$ has base
$\operatorname{ins}(b,f,j)$, satellite $s$, and visible count
$v+\boldsymbol{1}_{j<v}$. Its repair row at position $j$ has base
$\operatorname{rot}(b,j)s$, satellite $f$, and visible count $|b|+1$.
For a full row the completion path is its full transport path with new
letter $f$. For a short row at port $p<h-1$ the path is its one lift at
$p$. At port $h-1$ the path is, in order, the lift at $h-1$, the repair
rows at $h-1$ and $h$, and the lift at $h$. As before, concatenate paths
and propagate the source port. This defines the \emph{completed contour}
from each tagged initial port.

A \emph{circle} is a word of distinct letters, considered together with
all its rotations. Its displayed representative and its position in a
list of circles are retained for counting. For a contour with bases of
length $h+1$, the auxiliary circles have length $h$. We say that a circle
\emph{safely meets} a tagged completion port if either it contains $f$
and one of its rotations is the head of a row in the completed contour,
or the source contour is the full chart of a rotation of that circle.
The \emph{full chart} of a word $d$ and a letter $a$ outside $d$ consists
of the full rows with bases $\operatorname{rot}(d,j)a$, one for each
$0\le j<|d|$. This definition includes the circle $d$ used when a pending
class is closed. A list of circles is \emph{safe} for a contour if every
tagged initial completion port is safely met by at least one circle in
the list.

One useful sufficient witness for the first kind of meeting consists of
an actual source-row occurrence, the completion port propagated to that
occurrence, and a rotation identity
\[
 \operatorname{rot}(C,t)
 =\operatorname{ins}(\operatorname{head}(b,s,v),f,p).
\]
The first row of the local completion path has this head. Thus the
identity witnesses an actual vertex of the completed contour. Both the
occurrence and the propagated port matter; an equality of unlabeled
cyclic gaps is insufficient.

\begin{lemma}\label{lem-assembly-prefix}
Suppose a source path of port modulus $h$ starts at port $e$, where
$0\le e<h$. If its displacement is $S$, the displacement of its
transported path is
\[
 S+\left\lfloor\frac{e+S}{h}\right\rfloor.
\]
Immediately before the second transported row at a full source row whose
input port is zero, the transported prefix displacement is congruent to
$-e-1$ modulo $h+1$.
\end{lemma}

\begin{proof}
The formula is zero for the empty path. A full step subtracts one from
$S$. The floor decreases by one precisely when the old residue of $e+S$
is zero. Thus the claimed expression decreases by two at port zero and
by one at every other port. These are exactly the numbers of full rows
in the two possible full transport paths. A short step adds one to $S$.
The floor increases by one precisely at port $h-1$, so the expression
increases by two there and by one elsewhere. These are exactly the
displacements of the short transport paths. Induction proves the formula.
The floor is the Euclidean floor, also when $e+S$ is negative.

At a full port-zero occurrence, write $e+S=ht$. The first of its two
transported rows is full. Just before the second row, the displacement
is therefore $S+t-1=(h+1)t-e-1$, which gives the congruence.
\end{proof}

\begin{lemma}\label{lem-assembly-connectors}
In the setting of Lemma~\ref{lem-assembly-orbits}, suppose the source
contour has a consecutive run of $R$ full rows with bases
\[
 \operatorname{rot}(H,j)a,\qquad 0\le j<R,
\]
where $|H|=h$, $a$ is absent from $H$, and $R\le h$.
If $2g\lfloor R/g\rfloor\ge h$, the transported contours admit a safe
list of at most $2g^2$ circles. Each circle is supported on the source
alphabet together with $f$ and omits the satellite.
\end{lemma}

\begin{proof}
Put $L=\lfloor R/g\rfloor$, and let $S_0$ be the displacement before
the run begins. Fix a component with initial source port $e<g$, and
let $a_e$ be the residue of $e+S_0$ in $\{0,\ldots,g-1\}$.
The full source rows at phases
\[
 j=a_e+gu,\qquad 0\le u<L,
\]
all lie in the run. Indeed, $a_e+g(L-1)\le gL-1<R$.
Each is reached at port zero in exactly one of the $\ell$ laps. To
see this, its input port in lap $t$ is
$e+t\Delta+S_0-j$ modulo $h$. The congruence modulo $g$ holds by the
choice of $a_e$, and after division by $g$ the remaining congruence has
one solution for $t$ modulo $\ell$.

Write $H_e=\operatorname{rot}(H,a_e)$. For each $0\le t<g$, use the
two circles
\[
 C^-_t=\operatorname{ins}(H_e,f,g(L-1)+t),\qquad
 C^+_t=\operatorname{ins}(H_e,f,h-g+t).
\]
Their insertion positions lie between zero and $h-1$, because
$gL\le R\le h$. Both circles have length $h+1$, contain $f$, and
have distinct letters. At the second transported row of phase $a_e+gu$,
the head is $\operatorname{rot}(H_e,gu)$. In particular the transport
letter $z$ is absent from this head.

For $0\le j\le b<h$, splitting a word at positions $j$ and $b$ gives
\begin{align*}
 \operatorname{rot}(\operatorname{ins}(H_e,f,b),j)
   &=\operatorname{ins}(\operatorname{rot}(H_e,j),f,b-j),\\
 \operatorname{rot}(\operatorname{ins}(H_e,f,b),b+1)
   &=\operatorname{ins}(\operatorname{rot}(H_e,b),f,h).
\end{align*}
The second identity is a separate endpoint identity. It is needed because
the literal ports zero and $h$ of the advanced contour are distinct.

Given a port $gv+t$ with $0\le v<L$ and $0\le t<g$, choose
$u=L-1-v$ in the first identity for $C^-_t$. Thus the low circles meet
local ports $0,\ldots,gL-1$. The high circles, with $u$ varying through
$0,\ldots,L-1$, meet local ports $h-gL,\ldots,h-1$. For $C^-_0$ at
$u=L-1$, the second identity additionally meets port $h$.
The hypothesis $2gL\ge h$ says that the first two intervals cover
$0,\ldots,h-1$.

Finally fix a tagged initial completion port $J$ of this advanced
component. Lemma~\ref{lem-assembly-prefix} shows that its local port at
every selected second row is $J-e-1$ modulo $h+1$. Choose the local port
representative in $\{0,\ldots,h\}$ and use the corresponding circle
and occurrence above. The resulting literal identity proves safe
incidence from $J$. There are $2g$ circles for each of $g$ components,
so the total is $2g^2$. Coincident circles may remain as separate list
occurrences. Every circle uses only $H_e$ and $f$, which proves the
support statement.
\end{proof}

The literal hypothesis on the full run is necessary. Compatibility alone
allows a different full successor. For example, the three full bases
$012,120,201$ form a compatible contour, but their first two-letter
prefixes are not rotations of one fixed two-letter word. The positive-slot
construction must supply a run from an untouched normal batch, where
successive bases have precisely the form required in
Lemma~\ref{lem-assembly-connectors}.

\subsection{The finite universe above an old cyclic class}

Fix a word $d$ on $m\ge2$ distinct old letters and distinct new letters
$v_0,\ldots,v_{W-1}$. A cyclic class on these letters is \emph{above}
$d$ if deletion of all the $v_i$ leaves the cyclic class of $d$.
Define lists $U_i(d)$ recursively. Start with $U_0(d)=[d]$. For each
$b$ in $U_i(d)$, place
\[
 \operatorname{rot}(b,j)v_i,\qquad 0\le j<m+i,
\]
in $U_{i+1}(d)$, retaining the order of the parent list and the indicated
order of $j$. For example, from $d=ab$ and one new letter $x$ this gives
$U_1(d)=[abx,bax]$.

\begin{lemma}\label{lem-assembly-universe}
The list $U_W(d)$ contains exactly one representative of each cyclic
class above $d$. Its length is
\[
 D_{m,W}=\prod_{i=0}^{W-1}(m+i).
\]
If $B$ source-row occurrences have pairwise distinct base cyclic classes
above $d$, their complement in $U_W(d)$ is an actual list of length
$D_{m,W}-B$.
\end{lemma}

\begin{proof}
We prove the first assertion by induction on $i$. It holds for $i=0$.
Given a class above $d$ on the first $i+1$ new letters, delete $v_i$.
The remaining class has a unique representative $b$ in $U_i(d)$ by
induction. Rotate the larger word so that its unique occurrence of $v_i$
is last. The preceding word is a rotation of $b$, which gives existence
of a representative in the displayed list.

For uniqueness, equality of two larger cyclic classes implies equality
of their deleted parent classes. Thus their parent representatives agree.
Because $v_i$ occurs once and is last in both displayed representatives,
cyclic equivalence forces equality of those representatives as literal
words. A word of $m+i$ distinct letters has $m+i$ different rotations,
so the two rotation indices agree as well. Each parent therefore
contributes exactly $m+i$ new representatives. Multiplication over the
steps gives the formula for $D_{m,W}$.

Each selected source occurrence has a unique representative in $U_W(d)$.
The distinctness assumption makes these representatives distinct. Delete
these $B$ entries from the finite list. The resulting complement has
length $D_{m,W}-B$, as asserted. This counting operation does not change
the literal bases in the source contours.
\end{proof}

\subsection{Extension by an arbitrary number of letters}

We now specify the intermediate objects. Reserve distinct letters $s$ and
$f$ for the satellite and final completion. A \emph{stage of old size}
$m$ consists of an old alphabet $A_0$ of size $m$, a distinguished
ordinary letter $z$ outside $A_0$, a list of contours on $z\cup A_0$,
a list of auxiliary circles of length $m$, and a list of pending cyclic
classes on $A_0$. The alphabets omit $s$ and $f$. Each contour has
displacement divisible by $m$. Every circle has distinct letters and is
supported on $\{f,z\}\cup A_0$, and the circle list is safe for every
contour. Finally every cyclic class on $z\cup A_0$ either occurs as a
row base or is obtained by inserting $z$ into a pending cyclic class.
Write $R,Q,c,P$ for the row count, total charge, circle count, and pending
count. Counts refer to list occurrences.

We use the following existing transport facts. Transporting every
component from every port through a fresh ordinary letter preserves
closure, displacement divisibility and coverage. It multiplies row and
charge counts by $m+1$. Extending each circle through all its $m$ cyclic
insertion positions gives $mc$ circles, preserves distinct letters and
support, and safely meets every transported component. These statements
are the general one-letter level and circle transport results in the
supplied assembly library~\cite{halfar-assembly}. They apply to arbitrary
stages satisfying the preceding conditions.

\begin{lemma}\label{lem-assembly-width}
Let a stage have old size $m\ge2$, counts $(R,Q,c,P)$, old alphabet
$A_0$, and distinguished letter $v_0$. Let $W\ge1$, and choose distinct
fresh ordinary letters $v_1,\ldots,v_{W-1},z'$. Suppose that, above
every pending old class $d$, there is a family of closed source contours
on $A_0\cup\{v_0,\ldots,v_{W-1}\}$ with $B$ pairwise distinct source
classes above $d$ and total charge $q$. Suppose its transport through
$z'$ has displacement divisible by $m+W$ and admits a safe supported
list of $c_{\mathrm{adv}}$ circles of length $m+W$.
Then there is a stage of old size $m+W$ and distinguished letter $z'$
with counts
\begin{align*}
 R'&=A_{m,W}R+(m+W)PB,\\
 Q'&=A_{m,W}Q+(m+W)Pq,\\
 c'&=D_{m,W}c+Pc_{\mathrm{adv}},\\
 P'&=P(D_{m,W}-B),
\end{align*}
where $A_{m,W}=\prod_{i=1}^{W}(m+i)$ and $D_{m,W}$ is as in
Lemma~\ref{lem-assembly-universe}. The transport of the source family
is required to include every source-row and source-port pair once.
\end{lemma}

\begin{proof}
Set $A'_0=A_0\cup\{v_0,\ldots,v_{W-1}\}$. First transport the
existing contours successively through
$v_1,\ldots,v_{W-1},z'$, and extend their circles through the same
sequence. There are exactly $W$ transports: the old distinguished letter
$v_0$ already belongs to every existing component. At each step the
letter being added is absent from the current ordinary alphabet and from
both reserved roles. The one-letter results therefore apply inductively.
They give closed contours on $z'\cup A'_0$, displacement divisible by
$m+W$, valid supported circles, and safety from every completion port.
They multiply the existing row and charge counts by
$(m+1)\cdots(m+W)=A_{m,W}$ and the circle count by
$m\cdots(m+W-1)=D_{m,W}$.

For every old pending occurrence $d$, append its transported source
contours and its supplied circles. Their structural hypotheses put them
on the same alphabet with the same port modulus as the transported old
components. Taking the union preserves closure and displacement
divisibility component by component. It also preserves circle validity
and support. A circle which met a component before taking the union is
still in the union, so safety is preserved. Replace $d$ in the pending
list by the actual complement list from
Lemma~\ref{lem-assembly-universe}. Every new pending representative is
a permutation of $A'_0$.

It remains to prove the coverage condition. Let $y$ be an arbitrary
cyclic class on $z'\cup A'_0$. Delete the new letters in the reverse
order $z',v_{W-1},\ldots,v_1$. This produces a cyclic class $x$ on
$v_0\cup A_0$. Apply the old coverage condition to $x$.

If $x$ is the class of an existing row, reinsert
$v_1,\ldots,v_{W-1},z'$ in the gaps determined by $y$. The one-letter
coverage proposition at each transport gives a final transported row
in class $y$.

Otherwise $x$ is above an old pending class $d$. Delete $z'$ from $y$
and call the resulting source class $t$. Deleting all $W$ source
controls from $t$ gives $d$, so $t$ has one representative in $U_W(d)$.
If that representative was selected, the source family has a row with
base class $t$. The exhaustive transport inventory and the local
insertion coverage supply a transported row in the class $y$. If it was
not selected, its representative lies in the new pending complement,
and $y$ is obtained from that representative by inserting $z'$. These
two alternatives establish the new coverage condition.

Each pending old occurrence contributes $(m+W)B$ advanced rows and
charge $(m+W)q$, by the source-row and source-port inventory. It
contributes exactly $c_{\mathrm{adv}}$ circles and $D_{m,W}-B$ pending
classes. Adding these contributions to the transported old objects gives
all four displayed counts.
\end{proof}

For the source family of Proposition~\ref{prop-global-counts}, each
surviving candidate contributes the closed contour obtained by its
specified modifications. Order the finite set of surviving candidates
and list these contours in that order. The local distinctness theorem
and the exclusion of collisions between candidates show that all their
source-row cyclic classes are distinct. Their row and charge counts are
the sums of the counts of those same contours. In particular, selecting
this list requires no additional choice of rows or cyclic classes.

Each contour has an untouched normal batch and thus a full run of the
form required in Lemma~\ref{lem-assembly-connectors}. The verification
in Proposition~\ref{prop-global-counts} gives the inequality for its
length. Lemmas~\ref{lem-assembly-orbits}
and~\ref{lem-assembly-connectors} consequently provide the transported
contours, their complete source-port inventory, and their safe circle
lists. Thus the family satisfies every source hypothesis of
Lemma~\ref{lem-assembly-width}.

A family constructed on one fixed old order can be used above another
old order by a single bijection of the full finite alphabet. Apply this
bijection to the old letters, all source controls, the satellite, the
transport letter, and the completion letter together. Rotation,
insertion, prefixes, and equality of words commute with this relabeling.
It therefore preserves literal compatibility and every incidence
witness, as well as the counts. Relabeling only part of a frame would
not justify those equalities.

\subsection{An imported family for the final extensions}

For the last part of the schedule, we use an existing family with three
source controls. We state both its geometric properties and the numerical
bounds that will be used. Put
\[
 \mathcal G_a=\{g\in\mathbb Z:0\le g<a,\quad
   g\equiv1\pmod3,\quad 2g+5\le a\},
 \qquad
 \nu_a=
 \begin{cases}2,&a\text{ odd},\\8,&a\text{ even}.\end{cases}
\]

\begin{proposition}\label{prop-assembly-dense}
For every old cyclic order on $a\ge7$ letters, three distinct source
controls, a fresh transport letter, and two further distinct reserved
letters, there is a source family satisfying the hypotheses of
Lemma~\ref{lem-assembly-width} with width three. Its counts are
\[
 B_a=2|\mathcal G_a|a(a+12),\qquad
 q_a=16|\mathcal G_a|a,\qquad
 c_{\mathrm{adv},a}=(a+2)\nu_a.
\]
Writing $D_{a,3}=a(a+1)(a+2)$, these counts satisfy
$B_a\le D_{a,3}$ and $\nu_a\le8$. For $a\ge49$, they also satisfy
\[
 D_{a,3}\le4B_a,\qquad
 a(q_a+c_{\mathrm{adv},a})\le8B_a.
\]
\end{proposition}

This is the family \texttt{UniR.rFam} in the supplied assembly
library~\cite{halfar-assembly}. Its construction is valid for an
arbitrary full frame of distinct letters satisfying the two reserved-letter
exclusions. The declarations \texttt{rData\_ok} and
\texttt{rData\_safe} give literal source validity, distinct source classes,
and safe incidence from every required port. The constructor
\texttt{Uniform.famS} supplies the transported components, their
displacement divisibility, and their finite complement. Thus this
proposition gives actual families and circle lists for
Lemma~\ref{lem-assembly-width}.

The count formulas and the bounds $B_a\le D_{a,3}$, $\nu_a\le8$, and
$D_{a,3}\le4B_a$ are respectively \texttt{rowsR\_eq}, \texttt{chargeR\_eq},
\texttt{rowsR\_le}, \texttt{ncycR\_le}, and
\texttt{rows\_quarter} in \texttt{UniR/RRate.lean}.
The cost inequality is \texttt{cost\_le8} in
\texttt{UniR/RRate8.lean}. The count of advanced circles is
$(a+2)\nu_a$, since each of the original $\nu_a$ circles has $a+2$
extensions at transport. Whole-frame relabeling as described above
transfers these statements to the chosen old cyclic order and ordinary
letters. No parity restriction remains, and the family can extend an
already existing stage. It need not start a new construction.

\subsection{Initialization, padding, and the final word}

At old size $m$, choose no contours and no circles, and take every cyclic
class on the old alphabet to be pending. There are $(m-1)!$ such classes:
fix one old letter first and permute the remaining $m-1$ letters. Every
class on the old alphabet together with the distinguished letter lies
above one of them, by deletion of that letter. Thus this is a stage
with counts $(0,0,0,(m-1)!)$.

The identity
\[
 (m+W)D_{m,W}=mA_{m,W}
\]
and Lemma~\ref{lem-assembly-width} show that
\[
 R'+(m+W)P'=A_{m,W}(R+mP).
\]
Consequently every stage constructed from this initialization satisfies
$R+mP=m!$. Taking $W=1$ and an empty selected family gives a valid
one-letter extension with
\[
 (R',Q',c',P')=((m+1)R,(m+1)Q,mc,mP).
\]
Such extensions permit adjustment of the final alphabet size without a
congruence restriction.

At a final old size $h$, replace each pending representative $d$ by its
full chart with the current distinguished letter, and add the circle
$d$. The chart is closed, has $h$ rows, has zero charge, and has
displacement $-h$. It covers every class above $d$ and is safely met by
its own circle according to the definition. Thus the resulting closed
level has $h!$ rows, charge $Q$, and $c+P$ circles, and covers every base
cyclic class.

We use the following precise consequence of the imported word assembly
theorem~\cite{halfar-assembly}. Suppose $h\ge8$. A closed level on an ordinary alphabet
of size $h+1$, with $h!$ rows, charge $Q$, and a valid supported safe
list of $C$ circles, gives a covering word on $n=h+3$ letters of length
at most
\[
 F_3(n)+Q+hC+(a-2)C
 +(h+1-a)\min\{C,(n-1)^{\underline{h+1-a}}\},
 \qquad 2\le a\le h+1.
\]
The notation $x^{\underline{k}}$ means $x(x-1)\cdots(x-k+1)$.
The completed alphabet, consisting of the ordinary letters, satellite,
and completion letter, is relabeled bijectively to $\{0,\ldots,n-1\}$,
with the satellite labeled $9$ and the completion letter labeled $10$.
These labels are available because $n=h+3\ge11$. All other labels can
be assigned bijectively to the ordinary letters. The relabeling preserves
contiguous occurrences and meets the alphabet hypothesis of the imported
declaration \texttt{Closed.word}. The theorem supplies
an actual word containing every permutation, rather than an estimate
for an abstract count.

Taking $a=2$, $C=c+P$, and bounding the minimum by $C$ gives
\begin{equation}
 |w|\le F_3(n)+Q+(2h-1)(c+P).
 \label{eq-assembly-word}
\end{equation}
This is the completion estimate used in the finite transfer argument.
The source families of Proposition~\ref{prop-global-counts}, followed by
the families of Proposition~\ref{prop-assembly-dense}, maintain all of
these hypotheses through the extensions above.
